\section*{Summary}
Ember is a one-dimensional stellar-evolution code written to follow a
$0.1\Ms$ star from contraction, through hydrogen burning lasting several
trillion years, to a cooling helium white dwarf and its conditional fate if
nucleons decay.

A continuous calculation, started on the Hayashi track and run with one
program from deuterium burning onward, has passed $3.02\Tyr$ at 3121 K,
with the star still fully mixed and hydrogen at $X=0.2734$. This run uses a
bounded fixed-metal atmosphere. A second calculation selects the measured
metal- and hydrogen-dependent atmosphere families from the same Hayashi
starting prescription. Neither has yet reached the late burning shell.

A comparison calculation, started from a static hydrogen-burning
model on its own clock, includes gravitational settling of helium and
metals and reaches $4.002\Tyr$. Whole-star convection ends near
$3.55\Tyr$, central hydrogen falls below one part in a thousand near
$3.65\Tyr$, and luminosity and surface temperature pass shallow maxima
($0.006513\Ls$; 4754 K) before the model turns down at 4612 K, still
97\% powered by hydrogen burning. A companion calculation without settling
is about twice as luminous as the 1997 model of Laughlin, Bodenheimer and
Adams at the same mass and reaches each burning stage in about 60\% of the
time; MESA runs with grey atmospheres show the same sequence of events.

When the outer boundary of the settling model is switched to a published
pure-hydrogen white-dwarf atmosphere, a helium-3 shell flash develops. It
appears on every mass grid and mixing treatment tried, but every computed
history descends from one post-switch state, the shell was fading before
the switch, and no other calculation shows it. Its peak, its duration and
whether it occurs with the original boundary are not established; the
continuations through it are provisional, and the control with the
original atmosphere is pending.

The continuous settling calculation still needs atmosphere coverage through
cooling below 4600 K, coupled grain formation, hydrogen-free atmospheres and
a cold helium equation of state. Grey-atmosphere comparisons already reach
lower temperatures. The paper ends with a conditional timeline to the loss
of the remnant.
